\newcommand{\nomeEstado}{BAHIA}
\newcommand{\siglaEstado}{BA}
\newcommand{\nomeConcessionaria}{NEOENERGIA BAHIA}
\newcommand{\cnpjConcessionaria}{15.139.629/0001-94}
\newcommand{\termoCIP}{Sexto}
\newcommand{\presidenteConcessionaria}{Sr. Thiago Guth.}

\newcommand{\empresaResponsavel}{ABEL}
\newcommand{\nomeMunicipio}{Cachoeira}
\newcommand{\nomeMunicipioMaisculo}{\MakeUppercase{\nomeMunicipio}}
\newcommand{\IPestimada}{609351}
\newcommand{\cnpjMunicipio}{13.828.397/0001-56}
\newcommand{\sedePrefeitura}{Rua Ana Nery, nº 27, Centro (Centro Histórico), Cachoeira – BA, CEP 44300-000}
\newcommand{\nomePrefeito}{Eliana Gonzaga de Jesus}
\newcommand{\generoPrefeitura}{pela Prefeita}
\newcommand{\numeroContrato}{400/2025}
\newcommand{\quantidadeAditivos}{0}
\newcommand{\legitimidade}{C}
\newcommand{\publicacao}{Diário Oficial do Município}
\newcommand{\dataPublicacao}{12 de Dezembro de 2025}
\newcommand{\tituloClausula}{CLÁUSULA PRIMEIRA — OBJETO}
\newcommand{\clausula}{CLÁUSULA PRIMEIRA — OBJETO: 

O objeto do presente instrumento é a contratação de pessoa jurídica para prestação de serviços técnicos especializados, visando: Assessorar o Município na constituição de receitas de natureza tributária, inclusive habite-se de torres de geração eólica e solar, ISSQN de instituições financeiras e postos de atendimentos bancários, cartórios, construtoras e empresas optantes pelo Simples Nacional; Assessorar o Município na gestão, elaboração de auditorias e laudos técnicos, mediante a conferência das faturas de energia elétrica da Administração Direta e Indireta do Município, elaboração de memorial de cálculo de consumo e potência do parque de iluminação pública, apuração do modelo tarifário aplicado em cada unidade consumidora, assim como a verificação de possíveis isenções indevidas e/ou não repasse da Contribuição de Iluminação Pública (CIP), visando à repetição de indébitos decorrentes de cobranças indevidas (a maior) nas contas de energia elétrica de titularidade do Município de Cachoeira/BA. “Deverá a empresa efetuar levantamento dos créditos a que faz jus o Município, referentes aos pagamentos indevidos à concessionária de energia elétrica, conforme os prazos estabelecidos na Resolução Normativa da ANEEL nº 1.000, de 7 de dezembro de 2021, especialmente o art. 323, § 2º, inciso I, e suas alterações posteriores, podendo, para tanto, representar administrativamente o Município de Cachoeira/BA perante a Agência Nacional de Energia Elétrica — ANEEL, em todos os processos, reclamações, recursos e demais atos decorrentes de suas competências, para atender às necessidades do Município de Cachoeira/BA.”}
\newcommand{\telefone}{(75) 3425-1390}
\newcommand{\email}{ascomprefcachoeira@gmail.com}
